% Figures for this chapter: figures and tables/1_Introduction/
\chapter{Introduction}
\label{ch:introduction}

% Readable paragraph spacing within the Introduction only
\setlength{\parskip}{0.65em plus 0.15em minus 0.05em}
\setlength{\parindent}{0pt}

\section{Clinical Context and Motivation}
\label{sec:clinical-context}

\subsection{Anatomy of Coronary Arteries}
\label{subsec:anatomy}

The human heart is the central organ of the cardiovascular system and requires a continuous supply of oxygen and nutrients to function properly. This supply is provided by the coronary artery tree, a specialized vascular network that delivers blood directly to the heart muscle, known as the myocardium. The coronary arteries originate from the root of the aorta and surround the surface of the heart in a crown-like arrangement, from which the term \emph{coronary artery} is derived~\cite{clevelandclinic_coronary}. As illustrated in \Cref{fig:coronary-anatomy}, the coronary arterial system is mainly composed of two vessels: the right coronary artery (RCA) and the left coronary artery (LCA).

\begin{figure}[H]
  \centering
  \includegraphics[width=0.62\linewidth]{{Figure 1. Coronary Arteries.png}}
  \caption{Anterior view of the human coronary tree; right coronary artery (RCA) and left coronary artery (LCA) are highlighted in red, with other cardiac structures in blue. Adapted from Lynch et al.~\cite{lynch2010_coronary}.}
  \label{fig:coronary-anatomy}
\end{figure}

The RCA primarily supplies blood to the right atrium and right ventricle. In contrast, the LCA quickly divides into two major branches: the left anterior descending (LAD) artery and the circumflex (LCx), which together supply the majority of the left ventricle and the interventricular septum~\cite{clevelandclinic_coronary}. From these primary vessels, a network of smaller branches extends deep into the myocardium to ensure regional blood delivery across the entire heart.

A critical characteristic of the coronary circulation is its limited collateral connectivity. Unlike other vascular systems in the human body, coronary arteries generally lack sufficient alternative pathways to redirect blood flow when an obstruction occurs. Consequently, any narrowing or blockage in a proximal arterial segment directly compromises blood flow to all downstream regions supplied by that vessel~\cite{seiler2010}.

Additionally, coronary anatomy presents significant inter-patient variability in terms of vessel size, branching structure, and coronary dominance. Dominance is defined by which artery gives rise to the posterior descending artery (PDA) and can be classified as right-dominant, left-dominant, or codominant. Right-dominant circulation, in which the PDA originates from the RCA, is the most common configuration, followed by left dominance and codominance~\cite{wu2024}. Understanding this complex anatomical layout is fundamental, as it provides the basis for locating lesions, quantifying stenosis severity, and stratifying risk in coronary artery disease (CAD) diagnosis.

\subsection{Coronary Artery Disease}
\label{subsec:cad}

Coronary artery disease (CAD) remains one of the leading causes of morbidity and mortality worldwide~\cite{who2023_cvd}. It is characterized by the progressive accumulation of atherosclerotic plaque within the walls of the coronary arteries, leading to a narrowing of the vessel lumen known as stenosis. This narrowing restricts blood flow to the myocardium and, if left untreated, can result in myocardial ischemia or infarction~\cite{knuuti2020,mccullough2007}. \Cref{fig:cad-plaque} illustrates this pathological mechanism, contrasting a normal artery with one obstructed by plaque buildup.

\begin{figure}[H]
  \centering
  \includegraphics[width=0.62\linewidth]{{Figure 2. Coronary Artery Disease.png}}
  \caption{Coronary arteries narrowed by atherosclerotic plaque in coronary artery disease (CAD). Adapted from Kaiser Permanente~\cite{kaiser_atherosclerosis}.}
  \label{fig:cad-plaque}
\end{figure}

Early and accurate diagnosis is essential for effective risk stratification and patient management. In current clinical practice, coronary computed tomography angiography (CCTA) is the standard non-invasive imaging technique for CAD evaluation~\cite{knuuti2020}. CCTA provides detailed 3D reconstructions of the coronary tree, allowing clinicians to detect plaques, analyze their spatial location, and estimate the degree of luminal narrowing.

However, despite the high resolution of CCTA, the diagnostic workflow remains highly time-consuming and heavily dependent on manual and visual assessment. For instance, at Hospital de la Santa Creu i Sant Pau, reporting non-pathological or low-complexity cases requires approximately 30 minutes, while complex evaluations can take up to 2 hours~\cite{acebes2024}. This significant clinical burden drives the need for automated, data-driven tools capable of quantifying stenosis efficiently and supporting clinical decision-making, while fully preserving expert interpretability and clinical control.

\subsection{Stenosis Quantification: A Key Indicator of CAD Severity}
\label{subsec:stenosis-quantification}

Among the various indicators used to assess coronary artery disease, stenosis severity plays a central role in clinical diagnosis and risk stratification. Stenosis represents the reduction of the vessel lumen caused by atherosclerotic plaque accumulation, a condition directly associated with an increased risk of myocardial ischemia and adverse cardiac events~\cite{knuuti2020,garcia2014}. In clinical practice, this severity is quantified using geometric measurements such as the minimal lumen area (MLA) or percentage diameter stenosis (\%DS). These metrics evaluate the degree of narrowing by comparing the most constricted point of a lesion with a reference baseline, which is typically estimated from adjacent, relatively healthy vessel segments~\cite{garcia2014,hideo_kajita2019}.

Although these definitions are conceptually simple, accurate stenosis quantification presents significant clinical challenges. Coronary arteries exhibit complex, tortuous geometries, and diffuse atherosclerosis often makes it difficult to identify truly healthy reference regions. In addition, limitations in image resolution, segmentation errors, and noise in CCTA-derived data can highly influence measurements, leading to variability in stenosis estimates~\cite{kalisz2016,hideo_kajita2019}.

In routine clinical practice, stenosis assessment is often performed visually or using semi-automatic tools provided by proprietary software~\cite{acebes2024,cury2022}. Although these methods are widely used, they depend heavily on the clinician's interpretation and can vary between observers, which reduces consistency and reproducibility~\cite{hideo_kajita2019}.

Consequently, there is a strong clinical need to develop robust, transparent, and reproducible quantification methods. Automating these geometric measurements has the potential to standardize stenosis evaluation, reduce reporting time, and provide reliable quantitative data essential for downstream clinical tasks, such as CAD-RADS~2.0 assignment and patient prioritization.

\subsection{CAD-RADS 2.0 and the Need for Decision Support}
\label{subsec:cad-rads}

To standardize the reporting of coronary disease and facilitate clinical communication, the Coronary Artery Disease-Reporting and Data System (CAD-RADS~2.0) was established~\cite{cury2022}. This structured framework categorizes patients into distinct risk scores ranging from 0 to 5, primarily based on the most severe stenosis measured via CCTA. Each categorical score reflects the overall disease burden and provides specific clinical recommendations, offering a standardized pathway for subsequent patient management~\cite{cury2022}. Beyond maximal stenosis, CAD-RADS~2.0 also quantifies the overall extent of coronary plaque through the Segment Involvement Score (SIS), defined as the absolute number of assessable coronary segments exhibiting any measurable atherosclerotic plaque~\cite{cury2022}.

While CAD-RADS~2.0 significantly improves reporting consistency across institutions, its assignment still relies heavily on the manual interpretation of stenosis severity across multiple vessels and segments. This process requires clinicians to continuously synthesize quantitative geometric measurements with complex anatomical context. Consequently, evaluating these scans becomes highly cognitively demanding, particularly in time-constrained clinical environments~\cite{acebes2024}. In practice, manual scoring remains susceptible to inter-observer variability, especially in borderline cases or in patients presenting with multiple moderate lesions. Such discrepancies can lead to inconsistent diagnostic classifications, ultimately impacting downstream clinical decisions, further testing, and patient prioritization~\cite{acebes2024,hideo_kajita2019}.

To address these limitations, there is a growing clinical imperative to develop automated decision support tools. These solutions can span from rule-based algorithms mapping quantitative stenosis measurements to machine learning approaches that leverage structured imaging features. Crucially, these computational tools aim to provide objective, reproducible, and transparent suggestions to assist radiologists and cardiologists. By streamlining the CAD-RADS~2.0 assignment process, these systems reduce reporting times and variability while strictly preserving expert clinical judgment and final decision authority.

\section{Clinical Automation Challenges at Hospital de la Santa Creu i Sant Pau}
\label{sec:hsp-challenges}

\subsection{Current CAD Diagnostic Workflow and Associated Limitations}
\label{subsec:current-workflow}

The current diagnostic workflow for Coronary Artery Disease (CAD) at Hospital de la Santa Creu i Sant Pau involves a complex and coordinated series of steps, managed primarily by the specialized cardiac imaging unit (Unitat d'Imatge i Funció Cardíaca, UIFC). As illustrated in \Cref{fig:cad-workflow-hsp}, the clinical pathway begins when a patient presenting with symptoms such as chest pain is scheduled for a Coronary Computed Tomography Angiography (CCTA). Once the scan is acquired, the images are stored in the hospital's \emph{Picture Archiving and Communication System} (PACS) for further evaluation.

\begin{figure}[H]
  \centering
  \includegraphics[width=\linewidth]{{Figure 3. Current CAD Diagnostic Workflow HSP (2).pdf}}
  \caption{Current CAD diagnostic workflow at Hospital de Sant Pau, spanning coronary computed tomography angiography (CCTA) acquisition, picture archiving and communication system (PACS) storage, and radiology information system (RIS) reporting. The primary operational bottleneck (image assessment and reporting) is highlighted in red. Adapted from Ferrer Beltran~\cite{ferrer2024}.}
  \label{fig:cad-workflow-hsp}
\end{figure}

The core of this diagnostic pipeline relies heavily on the visual and manual analysis of these scans by expert cardiologists and radiologists to extract clinical metrics. These findings are then compiled into a clinical report, which the referring clinician evaluates to determine the appropriate subsequent medical actions, ranging from discharging a healthy patient to programming further invasive tests. While the overall infrastructure is well-established, the central phase of image analysis and reporting (highlighted in red in \Cref{fig:cad-workflow-hsp}) represents a critical operational bottleneck.

\subsection{The Bottleneck: Manual Image Analysis and Reporting}
\label{subsec:bottleneck}

Zooming into the specific image analysis and reporting phase (\Cref{fig:cad-workflow-hsp}), the limitations of the current clinical pathway become evident. Currently, clinicians rely on commercial software solutions, such as \emph{syngo.via}, to manually inspect the coronary arteries, differentiate segments, and calculate stenosis severity. Although these platforms offer semi-automated tools, the process remains highly dependent on the clinician's expertise to visually verify the vessels, identify bifurcations, and assess the plaque burden. Depending on the anatomical complexity and image quality, analyzing a single patient can consume around 30 minutes for healthy cases and 90 minutes for complex cases~\cite{acebes2024}.

Once the visual assessment is complete, the extracted metrics are manually entered into an internal custom-built platform called \emph{Filemaker}, which generates a template-based preliminary report. The specialist must then review, correct, and finalize this text before it is integrated into the \emph{Radiology Information System} (RIS).

\begin{figure}[H]
  \centering
  \includegraphics[width=\linewidth]{{Figure 4. Current Image Analysis and Reporting Workflow HSP (2).pdf}}
  \caption{Detailed representation of the CCTA image analysis and reporting workflow at Hospital de Sant Pau. Adapted from Acebes Pinilla~\cite{acebes2024}.}
  \label{fig:reporting-workflow-hsp}
\end{figure}

This heavy reliance on manual quantification and data entry presents significant operational challenges. Consequently, highly specialized experts spend considerable time performing repetitive, manual tasks on low-risk patients, rather than focusing their expertise on critical cases. This lack of an automated, prioritized workflow delays the diagnosis of severe cases and highlights a clear necessity for robust and automated solutions.

\section{State of the Art: Automation of CAD Diagnosis and Visualization}
\label{sec:state-of-art}

The integration of Artificial Intelligence (AI) and advanced computer vision in cardiac imaging has seen exponential growth over the past decade, driven by the need to optimize clinical workflows~\cite{litjens2017}. Deep learning algorithms, particularly Convolutional Neural Networks (CNNs), have demonstrated remarkable success in medical image analysis, encouraging a move from manual CCTA assessment to automated systems~\cite{slomka2017}. However, significant limitations remain regarding their practical implementation. For these new techniques to succeed, they must be practical for hospitals, integrate smoothly into current systems, and be easily adopted by the clinical staff~\cite{kelly2019}.

The prerequisite first step in any automated CAD assessment pipeline is the accurate segmentation and anatomical labeling of the coronary artery tree; before any downstream diagnostic metric can be calculated, the vessels must be correctly isolated and identified. Recent advancements have leveraged deep learning architectures, such as 3D U-Nets, to extract the vessel lumen from CCTA images~\cite{cicek2016}, while graph-based models like Graph Convolutional Networks (GCNs) are increasingly used to ensure topological consistency when labeling segments according to the SCCT-18 regional model~\cite{yang2020,leipsic2014}. Within Hospital de la Santa Creu i Sant Pau, preliminary research has explored these initial segmentation and labeling tasks using transformer-based models~\cite{clapers2025,sanchez2022}.

Following successful segmentation, extracting the precise geometric properties of the vessels is essential. The Vascular Modeling Toolkit (VMTK) stands out as one of the most prominent open-source libraries for medical geometry data extraction~\cite{izzo2018}. Specifically, algorithms implemented within VMTK are widely utilized to extract accurate vascular centerlines and compute cross-sectional areas~\cite{piccinelli2009,izzo2018}. These metrics provide the essential topological framework required to perform continuous Multi-Planar Reformations and subsequently evaluate vessel narrowing.

Once the 3D geometry of the coronary tree is established, the focus shifts to quantifying the disease burden. Current stenosis quantification techniques range from traditional intensity-based thresholding (which often struggles with severe calcifications) to modern machine learning models that predict stenosis severity directly from image features~\cite{zhang2024,hong2019}. Alternatively, geometric, non-ML approaches offer highly interpretable and robust methods by calculating the percentage of diameter or area stenosis based purely on the physical vessel profile, comparing the minimal lumen area against an interpolated healthy reference~\cite{hideo_kajita2019}. Recognizing the clinical value of algorithmic transparency, this project adopts that reference-window framework to quantify stenosis from extracted coronary geometry.

Beyond quantification, the effective visualization of these complex 3D metrics remains a critical challenge in the state of the art~\cite{borkin2011}. In the current clinical routine at Hospital de Sant Pau, specialists must manually navigate through cross-sectional slices and use standard viewing software to visually synthesize the data, assess the lesions, and manually fill out structured reports~\cite{acebes2024}. This cognitive load highlights a significant gap: the lack of dedicated, interactive visualization tools that bridge the gap between automated metrics and clinical reality. Modern research and clinical consensus emphasize that integrating 3D anatomical models alongside continuous luminal profiles and real medical images (such as Curved Multi-Planar Reformations) into a unified graphical interface is essential. Such integration allows clinicians to continuously validate automated findings against the ground-truth CCTA, drastically reducing the time and effort required to verify stenosis locations and seamlessly generate diagnostic reports~\cite{borkin2011}.

Several commercial software solutions have emerged to address these overarching needs, including platforms like \emph{Cleerly}~\cite{cleerly2024} and \emph{Artrya Salix}~\cite{artrya2024}, which use AI to analyze arterial plaque and assess ischemia risk from CCTA images. While these robust dashboards offer highly desirable visualization components, they often function as proprietary ``black boxes''. This lack of algorithmic transparency significantly hinders explainability, a feature strongly demanded by medical experts to verify how a specific CAD-RADS~2.0 category, plaque distribution, or stenosis percentage was computed. Finally, these standalone commercial platforms frequently struggle to integrate well into the highly specific IT infrastructures and patient prioritization workflows of public hospitals~\cite{kelly2019}.

\section{Project Scope and Contributions}
\label{sec:scope}

To address the clinical challenges outlined in \Cref{sec:hsp-challenges}, and to overcome the limitations of commercial tools discussed in \Cref{sec:state-of-art}, this project contributes directly to the comprehensive, transparent, and in-house AI-enabled diagnostic framework developed by the Dimension Lab at Hospital de la Santa Creu i Sant Pau~\cite{acebes2024}. Within the same hospital framework, prior student projects have already advanced both ends of the diagnostic pipeline. Upstream, coronary artery segmentation and labeling has been automated~\cite{clapers2025}; downstream, patient prioritization has also been addressed~\cite{ferrer2024}. Neither of these end stages falls within the scope of this project. The present work instead fills the gap between them, translating segmented coronary geometries into the quantitative findings required for clinical reporting.

Therefore, this bachelor's project focuses on automating the geometric quantification of coronary stenosis, assigning the associated CAD-RADS~2.0 category through rule-based scoring, and integrating these findings into an interactive clinical dashboard. Clinical optimization of disease quantification or CAD-RADS~2.0 scoring lies outside this scope. To realize this goal, this project develops a custom Python-based computational pipeline that implements the automated image analysis and feature extraction stages outlined in \Cref{fig:proposed-workflow}.

\begin{figure}[H]
  \centering
  \includegraphics[width=\linewidth]{{Figure 5. Proposed Image Analysis and Reporting Workflow HSP (2).pdf}}
  \caption{Contextual overview of the proposed automated CAD diagnostic framework at Hospital de la Santa Creu i Sant Pau, integrating CCTA-based analysis and Coronary Artery Disease Reporting and Data System (CAD-RADS~2.0) scoring. The contributions of this thesis (green dashed box) bridge the gap between existing hospital data systems (blue) and the final clinical report. Adapted from Acebes Pinilla~\cite{acebes2024}.}
  \label{fig:proposed-workflow}
\end{figure}

By seamlessly connecting these automated outputs with the hospital's reporting systems, this work aims to mitigate operational bottlenecks and reduce overall diagnostic times. However, it is crucial to emphasize that the solution developed in this project is intended solely as a clinical decision-support tool, not as a replacement for medical expertise. In the highly sensitive healthcare sector, human oversight remains essential. Rather than replacing detailed diagnostic software, the visualization interface is designed to act as a rapid initial assessment tool that presents an immediate, high-level overview of each patient's critical metrics. It assists specialists in quickly evaluating prioritized patient lists and identifying cases that require in-depth analysis. Thus, the automated findings are designed to streamline the workflow, ensuring that final diagnostic decisions and clinical validation remain in the hands of medical professionals.

\section{Objectives}
\label{sec:objectives}

The primary objective of this thesis is to design and develop an automated, end-to-end computational pipeline that bridges the critical gap between raw coronary segmentation and actionable clinical reporting. Rather than focusing on the standalone optimization of individual algorithmic components, the core contribution of this work is the functional integration of these stages into a unified, transparent workflow tailored to the clinical environment of Hospital de la Santa Creu i Sant Pau. This project is intended as a clinical decision-support framework rather than a diagnostic replacement; consequently, the project scope prioritizes workflow reliability, geometric consistency, and interface usability over achieving optimal clinical concordance against expert-level diagnoses.

To achieve this, the thesis aims to gain a deep understanding of the current clinical environment, contextualizing the pipeline within the hospital's diagnostic workflow. Building upon this, the project intends to implement computational geometry algorithms capable of accurately extracting cross-sectional areas and percentage area stenosis from 3D anatomical models, as well as to develop a standardized rule-based logic to translate these metrics into CAD-RADS~2.0 category assignments. Furthermore, this work aims to build an interactive visualization interface that synthesizes these automated metrics into an exploration session, facilitating clinician workflow, case triaging, and patient prioritization. The final objective is to conduct a validation framework that assesses pipeline robustness, geometric consistency, and overall interface usability, identifying the tool's current operational boundaries and future clinical potential.

% Restore default paragraph spacing for subsequent chapters
\setlength{\parskip}{0pt plus 1pt}
\setlength{\parindent}{15pt}
